\documentclass[11pt]{article}

\usepackage[margin=2.5cm]{geometry}
\usepackage{amsmath,amssymb,amsthm}
\usepackage{dsfont}
\usepackage{amsfonts}
\newcommand{\R}{\mathbb{R}}
\usepackage{bm}

\title{Exercise Sheet -- Calculus}
\author{}
\date{}

\begin{document}
\maketitle

This sheet is organized in two parts:
\begin{itemize}
  \item \textbf{Part A -- Fundamental exercises}
  \item \textbf{Part B -- Advanced and finance-oriented exercises}
\end{itemize}

The problems are designed to practice: gradients and Hessians, Jacobians, one- and multi-dimensional change of variables, polar coordinates, Gaussian integrals, integration by parts, divergence/Green identities, and related results introduced in the lecture, with some applications to finance (lognormal distributions, Black--Scholes pricing, Greeks such as Vega).

\bigskip
\hrule
\bigskip

\section*{Part A -- Fundamental Exercises}

\begin{enumerate}
% -----------------------------
\item \textbf{Gradients of scalar functions}

Compute the gradient of the following functions:

\begin{enumerate}
\item \( f(x,y)=x^2y+3y^2 \).
\item \( f(x,y,z)=e^{xy}+z^2 \).
\item \( f(\bm{x})=\|A\bm{x}\|^2 \) where \(A\) is a fixed matrix and \(\|\cdot\|\) is the Euclidean norm. Express \(\nabla f(\bm{x})\) in terms of \(A\).
\end{enumerate}

\bigskip
% -----------------------------
\item \textbf{Hessian and convexity}

Let
\[
f(x,y)=x^4+2x^2y^2+y^4.
\]

\begin{enumerate}
\item Compute the Hessian matrix \(H_f(x,y)\).
\item Determine whether \(f\) is convex on \(\mathbb{R}^2\). Justify your answer.
\end{enumerate}

\bigskip
% -----------------------------
\item \textbf{Jacobian of a vector-valued map}

Consider the transformation \(\Phi: \mathbb{R}^2\to\mathbb{R}^2\):
\[
\Phi(u,v)=(u\,e^{v},\,u^2+v).
\]

\begin{enumerate}
\item Compute the Jacobian matrix \(J_\Phi(u,v)\).
\item Compute \(\det J_\Phi(u,v)\).
\end{enumerate}

\bigskip
% -----------------------------
\item \textbf{One-dimensional change of variables}

Evaluate the following integral using an appropriate substitution:
\[
I=\int_0^1 \frac{1}{\sqrt{1-x^2}}\,dx.
\]

Specify clearly:
\begin{itemize}
\item the change of variable you choose,
\item the transformed bounds,
\item the final value of the integral.
\end{itemize}

\bigskip
% -----------------------------
\item \textbf{Polar coordinates and Jacobian}

Compute
\[
\int_{B(0,1)} (x^2+y^2)\,dx\,dy,
\]
where \(B(0,1)\) is the unit disk in \(\mathbb{R}^2\). Use polar coordinates and the determinant of the Jacobian of the transformation give the geometric interpretation.

\bigskip
% -----------------------------
\item \textbf{Multiple integrals and Fubini’s theorem}

Show, using Fubini’s Theorem, that
\[
\int_0^1\int_0^1 xy\,dx\,dy=\int_0^1\int_0^1 xy\,dy\,dx.
\]

Compute the common value.

\bigskip
% -----------------------------
\item \textbf{Multivariate chain rule}

Let \(g:\mathbb{R}^2\to\mathbb{R}\) be given by
\[
g(u,v)=u^2+v^2,
\]
and let \(\Phi:\mathbb{R}^2\to\mathbb{R}^2\) be given by
\[
\Phi(x,y) = (u(x,y),v(x,y)) = (x^2-y,\; xy).
\]
Consider the composition \(f=g\circ \Phi\), that is
\[
f(x,y)=g(\Phi(x,y)).
\]

\begin{enumerate}
\item Compute \(\nabla f(x,y)\) directly from the expression of \(f\).
\item Compute the Jacobian \(J_\Phi(x,y)\) and \(\nabla g(u,v)\).
\item Use the multivariate chain rule to compute \(\nabla f(x,y)\) via
\(\nabla f(x,y) = J_\Phi(x,y)^\top \nabla g(\Phi(x,y))\),
and verify it matches part (a).
\end{enumerate}

\bigskip
% -----------------------------
\item \textbf{Linear change of variables in \(\mathbb{R}^2\)}

Let \(D\) be the triangle with vertices \((0,0)\), \((1,0)\), and \((1,1)\), and consider the integral
\[
\int_D (x+y)\,dx\,dy.
\]

\begin{enumerate}
\item Define a linear transformation \((u,v)\mapsto(x,y)=A(u,v)+b\) that maps the unit square \([0,1]\times[0,1]\) to the triangle \(D\).
\item Compute the determinant of the Jacobian of the transformation.
\item Use this change of variables to compute the integral.
\item Check your answer by computing the integral directly (without change of variables).
\end{enumerate}

\end{enumerate}

\bigskip
\hrule
\bigskip

\section*{Part B -- Advanced and Finance-Oriented Exercises}

\begin{enumerate}
\setcounter{enumi}{10} % start numbering at 11

\item \textbf{Black--Scholes setup and lognormality}

Under the risk-neutral measure, a stock price satisfies
\[
 S_T = S_0\exp\left((r-\tfrac12\sigma^2)T+\sigma\sqrt{T}Z\right), \qquad Z\sim N(0,1).
\]

\begin{enumerate}
\item Using change of variables, derive the probability density function of \(S_T\).
\item Show that 
\(\mathbb{E}[S_T] :=
\int_{\R} 
S_0\exp\left((r-\tfrac12\sigma^2)T+\sigma\sqrt{T}z\right) \frac{e^{-z^2/2}}{\sqrt{2\pi}} 
\mathrm{d}z
=S_0e^{rT}\).
\end{enumerate}

\bigskip
% -----------------------------
\item \textbf{Call and put prices as integrals over the standard normal}

Let $S_T$ be defined by the previous exercise. 
Express the undiscounted European call price
\[
 C = \mathbb{E}[(S_T-K)^+]
 =
 \int_{\R} 
 \Big(
 S_0\exp\big( (r-\tfrac12\sigma^2)T+\sigma\sqrt{T}z\big) 
 -
 K
 \Big)^{+}
 \frac{e^{-z^2/2}}{\sqrt{2\pi}} 
\mathrm{d}z
\]
as a one-dimensional integral with respect to the standard normal density by changing variables from \(S_T\) to \(Z\).

\begin{enumerate}
\item Write \(C\) explicitly as an integral \(\int_{\mathbb{R}} \phi(z)\,\cdots\,dz\) where \(\phi\) is the standard normal density.
\item Write an analogous integral expression for the undiscounted European put price \(P = \mathbb{E}[(K-S_T)^+]\).
\end{enumerate}

\bigskip
% -----------------------------
\item \textbf{Vega of a European call option}

Let \(C(S_0,K,\sigma,T)\) be the Black--Scholes call price (undiscounted) and define
\[
\text{Vega} = \frac{\partial C}{\partial \sigma}.
\]

\begin{enumerate}
\item Starting from the integral representation in the previous exercise, differentiate under the integral sign with respect to \(\sigma\). Justify briefly why the differentiation and integration may be interchanged.
\item Perform appropriate changes of variables to simplify the derivative.
\item Show that Vega can be written in the standard Black--Scholes closed form
\[
\text{Vega} = S_0 \sqrt{T}\,\phi(d_1),
\]
for the appropriate \(d_1\).
\end{enumerate}

\bigskip
% -----------------------------
\item \textbf{Vega of a European put option and put--call parity}

Let \(P(S_0,K,\sigma,T)\) be the Black--Scholes put price (undiscounted) with the same parameters \((S_0,K,\sigma,T)\).

\begin{enumerate}
\item Repeat the previous exercise for the European put option and derive an expression for \(\displaystyle \frac{\partial P}{\partial \sigma}\).
\item Use put--call parity to show that the Vega of the put equals the Vega of the call with the same \((S_0,K,\sigma,T)\).
\end{enumerate}

\bigskip
% -----------------------------
\item \textbf{Digital option price as a probability}

Consider a European digital (cash-or-nothing) call option paying \(1\) at time \(T\) if \(S_T > K\), and \(0\) otherwise.

\begin{enumerate}
\item Express the undiscounted price of this digital option as an integral
\[
D = \mathbb{E}\big[ \mathbf{1}_{\{S_T > K\}} 
=
\int_{\R} 
\mathds{1}_{
 \Big(
 S_0\exp\big( (r-\tfrac12\sigma^2)T+\sigma\sqrt{T}z\big) 
 \geq
 K
 \Big)^{+} }
 \frac{e^{-z^2/2}}{\sqrt{2\pi}} 
\mathrm{d}z
.
\]
\item Use the lognormal model for \(S_T\) and a suitable change of variables to write \(D\) in terms of the standard normal cumulative distribution function \(\Phi\).
\item Compare your result with the derivative of the Black--Scholes call price with respect to the strike \(K\).
\end{enumerate}





\bigskip
% -----------------------------
\item \textbf{Greeks and Jacobians for a two-instrument portfolio}

Consider a portfolio consisting of two options with price vector
\[
 P(\theta)=\begin{pmatrix}C(\theta)\\[4pt] P(\theta)\end{pmatrix},\qquad
 \theta=(S_0,\sigma).
\]

\begin{enumerate}
\item Form the Jacobian matrix of partial derivatives of \(P\) with respect to \(\theta\), i.e.
\[
J_P(\theta)=
\begin{pmatrix}
\displaystyle \frac{\partial C}{\partial S_0} & \displaystyle \frac{\partial C}{\partial \sigma} \\[8pt]
\displaystyle \frac{\partial P}{\partial S_0} & \displaystyle \frac{\partial P}{\partial \sigma}
\end{pmatrix}.
\]
\item Interpret each entry of \(J_P(\theta)\) in terms of Greeks (Delta and Vega of each instrument).
\item Discuss how the determinant of \(J_P(\theta)\) relates to local invertibility of the mapping from parameters \((S_0,\sigma)\) to prices \((C,P)\). What does a small determinant indicate from a numerical/hedging perspective?
\end{enumerate}

\bigskip
% -----------------------------
\item \textbf{Laplacian, heat equation and Black--Scholes PDE}

The heat equation and Black--Scholes PDE are related by a change of variables.
The Black--Scholes PDE for a call option \(V(S,t)\) with risk-free rate \(r\) and volatility \(\sigma\) is
\[
\frac{\partial V}{\partial t}
+ \frac{1}{2}\sigma^2 S^2 \frac{\partial^2 V}{\partial S^2}
+ r S \frac{\partial V}{\partial S}
- r V = 0, \,\, V(S, T) = (S-K)^+
\]

\begin{enumerate}
\item Verify that the change of variable $
x = \ln\!\left(\frac{S}{K}\right), 
\qquad 
\tau = \frac{1}{2}\sigma^2 (T - t),
$
Describe (without full derivation) a change of variables that transforms the Black--Scholes PDE into the classical heat equation
\[
\frac{\partial u}{\partial \tau} = \frac{\partial^2 u}{\partial x^2}.
\]
Hint: use a logarithmic transform in \(S\) and a suitable time reversal.
\item Explain conceptually why Gaussian integrals appear in the closed-form pricing formulas for European options.
\end{enumerate}

\bigskip
% -----------------------------
\item \textbf{Stability of an ODE using Taylor expansion}

Consider the scalar ordinary differential equation
\[
y'(t)=f(y(t)),
\]
and let \(y^\ast\) be an equilibrium point, that is
\[
f(y^\ast)=0.
\]

In this exercise we study the stability of \(y^\ast\) using a Taylor expansion of \(f\).

\begin{enumerate}

\item \textbf{Translation of variables}

Define the perturbation variable
\[
z(t)=y(t)-y^\ast.
\]

\begin{enumerate}
\item Show that \(z(t)\) satisfies an ODE of the form
\[
z'(t)=g(z(t))
\]
and express \(g\) in terms of \(f\) and \(y^\ast\).
\end{enumerate}

\item \textbf{Taylor expansion near the equilibrium}

Assume that \(f\) is \(C^1\) in a neighbourhood of \(y^\ast\).  
Write the Taylor expansion of \(f\) around \(y^\ast\):
\[
f(y^\ast+z)= f(y^\ast)+ f'(y^\ast)z + R(z).
\]

\begin{enumerate}
\item Use that \(f(y^\ast)=0\) to simplify this expression.
\item Give an explicit form of the remainder term \(R(z)\) using the Lagrange form.
\end{enumerate}

\item \textbf{Linearized equation}

By neglecting the remainder term \(R(z)\) for small \(z\), derive the \textbf{linearized equation}
\[
z'(t)=\lambda z(t)
\qquad\text{with}\qquad \lambda=f'(y^\ast).
\]

\begin{enumerate}
\item Solve this linear equation explicitly.
\item For which sign of \(\lambda\) does \(z(t)\to 0\) as \(t\to+\infty\)?
\end{enumerate}

\item \textbf{Conclusion on stability}

\begin{enumerate}
\item State a condition on \(f'(y^\ast)\) that guarantees local asymptotic stability of \(y^\ast\).
\item State a condition on \(f'(y^\ast)\) that implies instability.
\end{enumerate}

\item \textbf{Application}

Consider the ODE
\[
y'(t)= y - y^3.
\]

\begin{enumerate}
\item Find all equilibrium points.
\item Compute \(f'(y^\ast)\) for each equilibrium.
\item Use the criterion above to classify each equilibrium as
\[
\text{stable, asymptotically stable, or unstable}.
\]

\end{enumerate}
\end{enumerate}
\end{enumerate}





\end{document}